\chapter{Requisiti}

\section{Requisiti funzionali}

I requisiti funzionali sono stati tradotti in funzionalita' concrete dei servizi e delle viste principali.

\begin{itemize}
\item Registrazione, login e accesso guest.
\item Consultazione del catalogo con filtri per categoria, disponibilita' e ricerca testuale.
\item Creazione e modifica di una configurazione CAD passo per passo.
\item Calcolo del design ammesso e dell'opzione successiva disponibile per ogni colonna.
\item Finalizzazione della configurazione con popolamento automatico del carrello.
\item Ricreazione del carrello da una configurazione finalizzata tramite riordino.
\item Gestione notifiche per variazioni di prezzo e disponibilita'.
\item Chatbot contestuale che usa catalogo e configurazione corrente per dare risposte pertinenti.
\item Consultazione degli ordini per l'utente e funzioni di gestione stato per l'amministratore.
\item Visualizzazione dei trend di vendita per la parte di reporting.
\end{itemize}

\section{Requisiti non funzionali}

Il progetto e' stato impostato con attenzione a modularita', chiarezza del codice e fruibilita' web:

\begin{itemize}
\item separazione netta tra frontend, gateway e microservizi backend;
\item uso di API REST semplici e coerenti con status code ed error model uniformi;
\item presenza di canali real-time solo dove il valore percepito e' reale, cioe' notifiche e chat;
\item persistenza per servizio, cosi' da non condividere database tra domini distinti;
\item controlli di accesso espliciti per admin e utenti autenticati;
\item usabilita' base in desktop e mobile, supportata da layout responsive nel frontend;
\item test di unita', integrazione ed e2e per ridurre regressioni sui flussi critici.
\end{itemize}

\section{Tracciabilita' con le note di analisi}

Le note presenti in \texttt{utilities} hanno guidato la classificazione iniziale dei requisiti, la mappa dei bounded context e la traccia di copertura test. Durante il controllo sul codice reale sono emersi alcuni aggiustamenti importanti:

\begin{itemize}
\item il frontend usa Pinia e non Vuex;
\item il reporting implementato e' focalizzato sul trend ordini, non su una suite ampia di report;
\item il cart service espone le operazioni effettive di carrello e checkout, mentre la popolazione da configurazione avviene internamente dal CAD;
\item il CAD implementato e' un motore di configurazione step-based con logica geometrica e BOM, non un editor collaborativo WebSocket nel codice corrente.
\end{itemize}
